% New actual consequences of the stipulated relative lifting theorem R2.
\subsection{A kernel of descent detected by the quadratic character}

\paragraph{LI.1 (Refuted): the proposed nilpotent-residual exclusion at four.}
Let $U_0=F'F^3$, and evaluate $F$ in
$\mathrm{SL}_2(\mathbb Z/4)$ by $x\mapsto S$, $y\mapsto U^{-1}$.
The image of $U_0$ is the whole group of order $48$, and its
nilpotent residual is its derived subgroup of order $12$.
In particular it contains $-U^2$. There is even an explicit
actual profinite word $w\in R_{\rm nil}(U_0)$ with this value.
Thus membership $f\in R_{\rm nil}(U_0)$ alone cannot force
the first descent character to vanish.

\emph{Proof.}
The finite group has abelianization $C_4$, so commutators together
with cubes generate it. Its lower central series has successive
orders $48,12,12$, as the complete finite calculation verifies.
For an explicit source word let $e_p\in\widehat{\mathbb Z}$ be
the idempotent with $p$-adic component one and every other prime
component zero. Put
\[
 A=(x^6y^6)^{e_2},\qquad B=(x^3y^{-3})^{e_3},\qquad
 w=[BAB^{-1},B].
\]
Both $A,B$ lie in $U_0$. The first is a pro-two element and the
second is pro-three. Their indicated commutator maps to one in
every finite nilpotent quotient of $U_0$, whose different Sylow
subgroups commute. Hence $w\in R_{\rm nil}(U_0)$.
At the displayed modulo-four evaluation, $A=-U^2$ and $B=R_0$,
with respective orders two and three, and
\[
 [R_0(-U^2)R_0^{-1},R_0]=-U^2.
\]
This is an actual source counterexample to the residual exclusion.
It is not claimed to be a gentle parameter. The finite computations
are in \texttt{check\_iterated\_lifting.py}.

\paragraph{LI.2 (Proved): a simultaneous full-kernel word with nontrivial frame residue.}
Let $\varphi:F\to\widehat B$ be the current map of LC.1, and let
\[
 K=\operatorname{core}_{\widehat B}(\ker\varphi)
   =\bigcap_{b\in\widehat B} b(\ker\varphi)b^{-1}
       \quad\text{as a closed subgroup of }F.
\]
Then the explicit word
\[
                   A_0=[[x^4,y^4],z^4]^2
\]
belongs to $F'\cap K$ and has $q_{12}(A_0)=[5I]$.

\emph{Proof.}
The kernel of $\varphi$ is normal in $F$, so the displayed core
requires only representatives for the six cosets $\widehat B/F$.
Their conjugations permute the three peripheral classes of $F$.
Under each resulting evaluation into $\widehat B$, one of the
three entries $x,y,z$ is conjugate to $s$ and has order two.
Its fourth power is therefore one. The iterated commutator
$[[x^4,y^4],z^4]$ vanishes if any of its three entries is one.
Thus $A_0$ lies in all six kernels in the full profinite target,
not just at a congruence level. It is a commutator word, so belongs
to $F'$ as well. For the original inclusion use
$x=U^2$, $y=L^2$. Direct multiplication modulo twelve gives
$q_{12}(A_0)=[5I]$; the accompanying program gives all seven
matrix values, the six other values being $I$.

\paragraph{LI.3 (Proved, using R2): the quadratic character is nontrivial on the kernel of descent.}
There exists an actual $k\in\mathcal H$ such that
\[
 \boxed{\ D(k)=1,\qquad\Delta(k)=1.\ }
\]
It can simultaneously be chosen with
\[
 f_k\in R_{\rm nil}(U_0),\qquad
 q_{12}(A_k)=[5I],\quad q_{12}(\beta_k)=1,\quad \rho_k=0,
 \qquad U_0=F'F^3.
\]
In particular it is invisible on every absolute pro-prime pure action.
This is a full profinite existence theorem, not a finite survival result.

\emph{Proof.}
Let $J\triangleleft\widehat B$ be any open normal subgroup, and set
\[
 N_J=U_0\cap\operatorname{core}_{\widehat B}
                              \bigl(\varphi^{-1}(J)\bigr).
\]
This is open and normal in $\widehat B$, contained in $F$.
The word $A_0$ of LI.2 belongs to every $N_J$.
Retain $Q=\mathrm{SL}_2(\mathbb Z/12)/\{\pm I\}$ and
$q=q_{12}:\widehat B\to Q$. R2's finite calculation gives
$q(U_0)=V_4\times Q_8$, a two-group. Consequently $q(N_J)$
is a two-group containing $q(A_0)=[5I]$, and $q$ factors through
$\Gamma_2=\widehat B/O^2(N_J)$.

Put $D_J=\widehat B/N_J$ and retain the actual joint image
\[
 E_J=\operatorname{im}(\widehat B\to Q\times D_J).
\]
Its kernel over $D_J$ is exactly $q(N_J)$, and the kernel of
$\Gamma_2\twoheadrightarrow E_J$ is pro-two. The element
$([5I],1)$ belongs to $E_J$, because it is the image of $A_0$,
and it centralizes the image of $r$. The order of that image is
three, detected by the retained $S_3$ quotient. Coprime
centralizer lifting therefore produces
\[
 c_2\in C_{\Gamma_2}(r),\qquad
 c_2\longmapsto([5I],1)\in E_J.
\]
Here the elementary lifting argument is the one used in R2:
two complements to a pro-two kernel for the same order-three
subgroup are conjugate by the kernel, with the same generator
image; finite complement conjugacy and compactness give the
profinite assertion.

Use the identity automorphism in every relative pro-$p$ component
of $\widehat B/R_{\rm nil}(N_J)$. Give $s,\tau$ conjugator one,
and give $r$ conjugator $c_2$ at two and one at all other primes.
All these coordinates have the same value one in $D_J$.
The relative product theorem and R2 lift them to an actual full
peripheral automorphism. Normalize its cusp by the inverse of
the lifted cusp conjugator, which lies in $R_{\rm nil}(N_J)$.
The resulting normalized gentle pair $k_J=(1,f_J)$ has
\[
 f_J,\beta_J\in R_{\rm nil}(N_J),\qquad
 q(A_J)=[5I],\quad q(\beta_J)=1.
\]
Since $R_{\rm nil}(N_J)\subseteq N_J'\subseteq F'$, its
harmonic frame has abelianization zero, hence $\rho_J=0$.
R4 now gives $\Delta(k_J)=1$. Moreover
$\varphi(f_J)\in J$.

These are witnesses for a compactness argument on the actual
compact group $\mathcal H$. Impose the closed conditions
\[
 f\in R_{\rm nil}(U_0),\quad\beta\in R_{\rm nil}(U_0),
 \quad q(A)=[5I],\quad q(\beta)=1,
 \quad\varphi(f)\in J.
\]
Every finite collection is met by the construction for the
intersection of its open normal subgroups $J$.
Indeed $N_J\subseteq U_0$ implies
$R_{\rm nil}(N_J)\subseteq R_{\rm nil}(U_0)$.
The normalized parameter and the unique pure frames depend
continuously on the gentle element, so all these sets are closed.
Compactness gives an actual $k$ satisfying every condition and
$\varphi(f_k)=1$.

In LC.1 this says $C=1$, $p=q=a=0$, and $\beta_D=1$.
The normalized descendant fixes both $s$ and $\tau$, hence is
the identity. Thus $D(k)=1$, while the retained frame values
give $\Delta(k)=1$. The inclusion
$R_{\rm nil}(U_0)\subseteq R_{\rm nil}(F)$ proves the final
absolute pronilpotent invisibility assertion.

\paragraph{LI.4 (Proved): consequences for fixed subgroups and iterated characters.}
The descent endomorphism $D$ is not injective. Both
$\operatorname{Fix}(D)$ and the correctly marked
$\operatorname{Fix}(\operatorname{Ad}_c\circ D)$ are proper
closed subgroups of $\widehat{\mathrm{GT}}_0$. For every $n\geq1$,
\[
 \Delta\circ D^n\ne\Delta,
\]
and $\Delta$ does not belong to the finite linear span of these
positive iterates. More strongly,
\[
 (D,\Delta):\widehat{\mathrm{GT}}_0
       \longrightarrow\operatorname{im}(D)\times\mathbb F_2
\]
is onto; the corresponding multiplier-one statement is also true.
Using the independently proved arithmetic marking comparison
$D(\sigma)=c\sigma c$ on $G_{\mathbb Q}$, every
$\chi_n=\Delta\circ D^n$ vanishes on the Galois image.
Consequently $\chi_1$ is either zero or linearly independent of
$\Delta$ as a relative quadratic character. Its nonvanishing
is not established by the present construction.
Moreover, neither fixed subgroup contains $\ker\Delta$.
In particular the proposed equality
$\operatorname{Fix}(\operatorname{Ad}_c\circ D)=\ker\Delta$
is false even in the corrected original marking. More precisely,
$\ker D\cap\mathcal H$ contains a closed subgroup
$Z\simeq\mathbb Z_2$ on which $\Delta$ is reduction modulo two.
Every nonidentity element of $2Z$ lies in $\ker\Delta$ and in
neither fixed subgroup. No pentagon assertion about these
elements is made.

\emph{Proof.}
Use the actual element $k$ of LI.3. It is nonidentity since
$\Delta(k)=1$, but both descent endomorphisms send it to one.
Thus it is in neither fixed subgroup. Every positive iterate
$D^n(k)$ is one. Evaluating the proposed equality of characters,
or a finite linear combination of positive iterates, at $k$
proves the stated distinctions. Finally, for any chosen preimage
$g$ of a point of $\operatorname{im}(D)$, the elements $g,gk$
have the same descendant and opposite $\Delta$ values.
Both stay in $\mathcal H$ when $g\in\mathcal H$.
For the arithmetic assertion use the simultaneous positive-loop
comparison proved in A4.M3: $D$ preserves the Galois image and
acts there by conjugation by $c$. R4 therefore makes every
$\chi_n$ vanish on that image. Since $\chi_1\ne\Delta$,
the only remaining possibilities in a vector space of quadratic
characters are $\chi_1=0$ or linear independence.

For the fixed-subgroup comparison, R7 implies $k^2\ne1$,
because $k\in\mathcal H$ is nonidentity and $\mathcal H$ is
torsion-free. But $D(k^2)=1$ and $\Delta(k^2)=0$.
Both $D$ and $\operatorname{Ad}_c\circ D$ send $k^2$ to one,
so neither fixes it. For the stronger closed-subgroup statement
put $k_2=k^{e_2}$, using the prime idempotent of LI.1.
Continuity gives $D(k_2)=1$ and $\Delta(k_2)=1$.
The closed subgroup $Z=\overline{\langle k_2\rangle}$ is
procyclic and pro-two, nontrivial by its quadratic value, and
torsion-free by R7. Hence $Z\simeq\mathbb Z_2$.
Its nonzero homomorphism to $\mathbb F_2$ is reduction modulo
two, with kernel $2Z$. The two descent endomorphisms are
trivial on $Z$, proving all remaining assertions.

\paragraph{LI.5 (Proved, using R2): an actual independent affine-pentagon obstruction.}
Let $p,t$ generate $S_3$ with $p^3=t^2=(pt)^2=1$, and define
$a_f\in\mathbb F_3$ by $f(p,t)=p^{a_f}$.
There exists an actual multiplier-one gentle pair with
\[
                  a_f=0,\qquad \rho=-1\pmod3.
\]
It is identity on every absolute pro-prime pure action. Combined
with the independently verified five-point affine pentagon formula
$(a_f-\rho)e_1$, this is an actual profinite counterexample to the
pentagon, detected by that genuine five-point quotient.

\emph{Proof.}
Let $N$ be the $\widehat B$-normal core of the kernel of
$F\to S_3$, $x\mapsto p$, $y\mapsto t$. Every geometric
permutation of the boundary triple gives a generating $S_3$
triple with orders $3,2,2$. The two explicit words
\[
 b_x=x^2yx^2y^{-1},\qquad b_y=y^2xy^2x^{-1}
\]
vanish under every such evaluation. For $b_x$, either the image
of $x$ has order two, or it has order three and the image of $y$
is an involution inverting it. These are exactly the two cases
and both give one. The proof for $b_y$ is identical. Their
abelianizations are $(4,0)$ and $(0,4)$, so $N$ maps onto
$F/F'F^3=\mathbb F_3^2$, and $b_0=b_xb_y$ has image $(1,1)$.

Use the actual finite quotient $Q_3=\widehat B/(F'F^3)$ of
order $54$. Its pure kernel is $\mathbb F_3^2$, and $s$
interchanges the two coordinates. The image of $N$ is this
three-group, and the image of $b_0$ is the nonzero diagonal
translation, which centralizes $s$.
In the actual joint quotient of $Q_3$ and $\widehat B/N$,
the element $((1,1),1)$ therefore lifts to a centralizer of $s$
in $\widehat B/O^3(N)$, because $s$ has order two and the
kernel is pro-three.

Choose the identity relative automorphism, this conjugator
for $s$ at three, and conjugator one for $s$ at every other
prime and for $r,\tau$ at every prime. R2 and cusp normalization
give an actual full gentle pair with
\[
 f\in R_{\rm nil}(N)\subseteq N,\qquad
 \beta^{\rm ab}\equiv(1,1)\pmod3.
\]
The correction at the cusp lies in $R_{\rm nil}(N)$ and thus
has zero abelianization, so the retained frame congruence is
literal. It gives $\rho=-1\pmod3$, whereas $f\in N$ gives
$f(p,t)=1$ and $a_f=0$. Finally
$R_{\rm nil}(N)\subseteq R_{\rm nil}(F)$ proves the claimed
absolute pronilpotent invisibility. The finite seed identities
are checked on all eighteen generating pairs of $S_3$ by the
accompanying program. Existence of the full pair is supplied
by the displayed relative lift, not inferred from that enumeration.

\paragraph{LI.6 (Proved): all iterated quadratic tests miss an actual nonpentagonal element.}
There is an actual pro-three element $h\in\mathcal H$ which is
identity on every absolute pro-prime pure action, fails the pentagon,
and satisfies
\[
                   \Delta(D^n h)=0\qquad(n\geq0).
\]
Its closed procyclic subgroup is isomorphic to $\mathbb Z_3$.
Thus the intersection of all these quadratic kernels, even together
with all absolute pronilpotent tests, is strictly larger than its
intersection with $\widehat{\mathrm{GT}}$.

\emph{Proof.}
Take the actual element $h_0$ from LI.5 and put
$h=h_0^{e_3}$, where $e_3$ is the prime idempotent of LI.1.
The two maps $a_f,\rho\bmod3$ are additive on $\mathcal H$.
For $a_f$, evaluate the normalized action in the preserved quotient
$S_3=\langle p,t\rangle$: it fixes $p$ and conjugates $t$
by a power of $p$, so its action on $\langle p\rangle$ is trivial.
The pair multiplication law therefore adds the exponent $a_f$.
Additivity of $\rho$ is R4's frame multiplication law.
It follows that $(a_f-\rho)(h)=1$, since $e_3\equiv1\pmod3$.
The affine pentagon detector consequently remains nonzero.
Absolute pro-prime invisibility is preserved by taking a profinite
power. The closed subgroup generated by $h$ is pro-three, so every
continuous homomorphism from it to $\mathbb F_2$ is zero.
Apply this to each $\Delta\circ D^n$.
Finally $h\ne1$, and R7 excludes nontrivial finite-order elements
in $\mathcal H$. A nontrivial torsion-free procyclic pro-three
group is $\mathbb Z_3$, which proves the last assertion.
This conclusion only identifies an actual nonpentagonal element
in the joint zero locus. It assumes no vanishing of any
$\chi_n$ on all of $\widehat{\mathrm{GT}}$.

\paragraph{LI.7 (Proved): the cubic pentagon character measures a descent discrepancy.}
On $\mathcal H$ let $\psi(\Phi)=a_f-\rho\pmod3$, with
$a_f$ as in LI.5. In the prompt's marking the following are
identities of continuous homomorphisms on $\mathcal H$:
\[
 \rho\circ D=-a_f\pmod3,
 \qquad
 \boxed{\ \psi=-\rho\circ D-\rho\pmod3.\ }
\]
Consequently $\rho\circ D=-\rho\pmod3$ is a necessary
pentagon equation, and LI.5 gives an actual gentle pair violating
it. These identities concern $\mathcal H$; they do not assert
a full-multiplier extension of $\psi$.

\emph{Proof.}
Use the actual notation of LC.1. Reducing
$H=\varphi(f(z,x))$ modulo the pure subgroup gives
$f(r^{-1},s)=r^{-a_f}$ in $S_3$. Here simultaneous conjugation
shows that the defining exponent $a_f$ is independent of the
chosen generating pair $(p,t)$ of orders three and two, and
replacing $p$ by $p^{-1}$ changes the displayed power in the
specified way. Thus $l=-a_f\pmod3$.
LC.2 proves $a=l\pmod3$, whereas LC.1 gives
$a=q-p$, $p+q=0\pmod3$ and $\rho_D=-q$.
It follows that $a=-q=-a_f\pmod3$, hence
$\rho_D=-a_f$. Subtracting $\rho$ gives the boxed identity.
The genuine five-point affine pentagon calculation makes
$\psi=0$ necessary for the pentagon, and LI.5 proves that this
necessary equation is not a consequence of the gentle relations.

\paragraph{LI.8 (Proved, using R2): the character survives beyond the third derived action.}
The element $k$ of LI.3 can be chosen so that, in addition,
\[
       f_k,\beta_k\in R_{\rm nil}(F'')\subseteq F''',
       \qquad A_k\in F'',
       \qquad D(k)=1,\qquad\Delta(k)=1.
\]
Here $F''$ is the closed second derived subgroup of the full
free profinite group, and $F'''=(F'')'$ is its closed third
derived subgroup. The genuine normalized action of $k$ is
identity on $F/F'''$, hence on $F/F''$, on the completed Alexander module
$F'/F''$, and on $K/K'$ for every open subgroup $K$ containing
$F''$. In particular the canonical homology of finite abelian
covers, including the Fermat covers, does not contain a line
affording $(-1)^\Delta$ over a field of characteristic different
from two. The same is true of the completed Alexander module
after extension to such coefficients, and of the homology of
any finite metabelian cover preserved by the gentle action.
This assertion concerns these canonical modules; it does not
assert the analogous statement for arbitrary nonabelian covers
such as the full level-twelve cover.
It is also identity on $\widehat B/R_{\rm nil}(F'')$, hence
on $\widehat B/\widehat B'''$. In particular it is invisible
in every finite modular-source quotient in which the image
of $F''$ is nilpotent.
Consequently no action on any finite solvable quotient of
$\widehat B$ of derived length at most three can detect
$\Delta$, even if all such actions are retained jointly.
Together with the independently proved finite characteristic
$\widehat B$ detector of derived length four in FP.35, this
makes four the exact minimum derived length for such a
finite modular-source action detector. This is a statement
about quotients of $\widehat B$, not about the derived length
of arbitrary finite representations of the gentle group.

\emph{Proof.}
Replace the seed of LI.2 by the word
\[
                    A_1=[[x^4,y^4],[y^4,z^4]].
\]
Both inner commutators belong to $F'$, so $A_1\in F''$.
Under every $\widehat B/F$ conjugate of $\varphi$, one of
$x^4,y^4,z^4$ is one. If it is $x^4$ or $z^4$, the respective
inner commutator is one; if it is $y^4$, both are one.
Consequently $A_1\in\operatorname{core}_{\widehat B}\ker\varphi$.
Direct multiplication in the original signed level-twelve
representation gives $q_{12}(A_1)=5I$. The complete six
orbifold values and the original value are checked in
\texttt{check\_iterated\_lifting.py}.

Let $J\triangleleft\widehat B$ be open normal, and let
$M\triangleleft\widehat B$ be open normal with $F''\subseteq M$.
Set
\[
 N_{J,M}=U_0\cap M\cap
       \operatorname{core}_{\widehat B}(\varphi^{-1}J).
\]
The seed $A_1$ belongs to every $N_{J,M}$. Hence its level-twelve
image is a two-group containing $[5I]$, exactly as in LI.3.
The same coprime lift at $r$, followed by R2 and cusp
normalization, gives an actual $k_{J,M}\in\mathcal H$ with
\[
 \begin{gathered}
 f_{J,M},\beta_{J,M}\in R_{\rm nil}(N_{J,M}),
 \qquad A_{J,M}\in N_{J,M},\\
 q_{12}(A_{J,M})=[5I],\quad q_{12}(\beta_{J,M})=1,
 \quad\rho_{J,M}=0,\quad\varphi(f_{J,M})\in J.
 \end{gathered}
\]
The additional statement about $A_{J,M}$ follows from its
prescribed base coordinate one in $\widehat B/N_{J,M}$.
The cusp correction lies in $R_{\rm nil}(N_{J,M})$ and
therefore preserves membership of all three entries in
$N_{J,M}$. It also preserves their retained finite values.
Thus $f_{J,M},\beta_{J,M},A_{J,M}\in M$.

In fact $R_{\rm nil}(N_{J,M})\subseteq R_{\rm nil}(M)$:
every finite nilpotent quotient of $M$ restricts to one of
$N_{J,M}$. Impose the closed conditions of LI.3 together with
$f,\beta\in R_{\rm nil}(M)$ and $A\in M$, for every such $M$.
A finite collection is
met by taking the intersections of its $J$'s and $M$'s.
The continuity of the normalized parameter and pure frames
and compactness of $\mathcal H$ give a simultaneous limit.
For completeness the two exact closed-subgroup identities are
\[
 \bigcap_{\substack{M\triangleleft_o\widehat B\\F''\subseteq M}}M=F'',
 \qquad
 \bigcap_{\substack{M\triangleleft_o\widehat B\\F''\subseteq M}}
          R_{\rm nil}(M)=R_{\rm nil}(F'').
\]
The first is the separation of the profinite quotient
$\widehat B/F''$ by its finite quotients. For the second,
monotonicity gives containment from right to left. Conversely
let $q:\widehat B\to Q$ be any finite quotient killing
$R_{\rm nil}(F'')$. The finite group $q(F'')$ is nilpotent,
since $F''/R_{\rm nil}(F'')$ is pronilpotent.
Take $M=F''\ker q$, which is open normal and contains $F''$.
Then $q(M)=q(F'')$ is nilpotent, so $q(R_{\rm nil}(M))=1$.
All such $q$ separate the quotient by $R_{\rm nil}(F'')$,
proving the reverse containment. Normality here is legitimate:
$F''$ is characteristic in the normal subgroup $F$, and its
nilpotent residual is characteristic in $F''$.

The limit therefore has $f,\beta\in R_{\rm nil}(F'')\subseteq F'''$
and $A\in F''$, and retains
$\varphi(f)=1$, the two frame values, and $\rho=0$.
LC.1 and R4 again give $D(k)=1$ and $\Delta(k)=1$.
It also retains the absolute pronilpotent invisibility of LI.3.

The normalized genuine automorphism fixes $x$ and sends
$y$ to $f_k^{-1}yf_k$. Since $f_k\in F'''$, it induces identity
on $F/F'''$. Its restriction to $F'/F''$ is therefore identity.
If $K$ is open and contains $F''$, then $F'''\subseteq K'$.
Moreover $k$ preserves $K$ since it is identity on $F/F'''$.
For every $w\in K$ the element $k(w)w^{-1}$ lies in
$F'''\subseteq K'$, so the action on $K/K'$ is identity.
These are actual actions from the fixed normalized pair, so
no inner-automorphism ambiguity enters the assertion.
The normalized modular extension has
$k(s)=sf_k$, $k(r)=f_k^{-1}r$ and $k(\tau)=\tau$.
Since $f_k\in R_{\rm nil}(F'')$, it is identity on
$\widehat B/R_{\rm nil}(F'')$. Any finite quotient whose
image of $F''$ is nilpotent kills this residual and thus sees
the identity action. The inclusion
$R_{\rm nil}(F'')\subseteq F'''\subseteq\widehat B'''$
gives the asserted identity on $\widehat B/\widehat B'''$.
Every finite solvable quotient
of $\widehat B$ of derived length at most three kills
$\widehat B'''$; hence $k$ preserves each such quotient map
and acts identically on its target. Its nonzero $\Delta$
therefore rules out every detector from that joint collection.
FP.35 supplies the matching upper bound by an actual finite
characteristic quotient with a faithful enough outer-action
detector; no statement about the full pentagon is used here.
When $K$ contains $F'$, it is preserved by the full gentle
group as well: its image in $F^{\rm ab}$ is an open additive subgroup,
preserved by multiplication by every unit
$\lambda\in\widehat{\mathbb Z}^{\times}$.
For general metabelian covers we impose this stability when
regarding their individual homology as a module for the full
group; characteristic covers in particular satisfy it.

Over a field of characteristic different from two, a line
affording $(-1)^\Delta$ would make this same element $k$
act by $-1$. Since its action on each indicated module is
identity, no such line exists. More generally, for any nonzero
one of these coefficient modules $V$, its twist
$V\otimes(-1)^\Delta$ is not isomorphic to $V$ as a
$\widehat{\mathrm{GT}}_0$ module: an intertwiner would satisfy
$T=-T$ by applying $k$ and hence be zero. Their restrictions
to $G_{\mathbb Q}$ are canonically identical because R4
gives $\Delta|_{G_{\mathbb Q}}=0$.

Finally let $V_*$ denote the actual joint image of
$\widehat{\mathrm{GT}}_0$ in its absolute pro-prime pure
actions, its normalized full action on $F/F'''$, and the
canonical homologies of all finite abelian covers and all
finite metabelian covers preserved by the full group; include
also the canonical outer action on
$\widehat B/R_{\rm nil}(F'')$ (and hence on
$\widehat B/\widehat B'''$).
The pure and homology actions are the continuous genuine
actions from normalized pairs, while the last action is
explicitly an outer action. Cover stability has been
specified above. The retained closed condition
$f_k\in R_{\rm nil}(U_0)$ gives the first invisibility,
and the preceding proof gives all the other ones. Thus
\[
 \widehat{\mathrm{GT}}_0\longrightarrow V_*\times\mathbb F_2,
 \qquad g\longmapsto(\operatorname{image}(g),\Delta(g))
\]
is onto: $g$ and $gk$ have identical first coordinate and
opposite second coordinates. In particular $\Delta$ is not
a function of this entire joint collection of action and
homology data. This conclusion does not concern arbitrary
solvable quotients.

\paragraph{LI.9 (Proved): one actual witness passes the entire adelic primitive-cube-root Burau pentagon family.}
Let $\mathcal O=\mathbb Z[\omega]/(\omega^2+\omega+1)$,
and use the three literal Burau matrices $A,B,C$ of FP.36.2.
For every prime $p$ and $n\geq1$, reduce their projective
classes modulo $p^n$. This gives a genuine finite quotient
$\eta_{p^n}$ of $L=B_4/\langle z_4\rangle$, and hence an
induced genuine five-point quotient as in FP.36.
There exists one actual $k=(1,f_k)\in\mathcal H$ such that
\[
 \Delta(k)=1,\qquad Dk=1,\qquad\rho(k)=0,\qquad
 f_k,\beta_k\in R_{\rm nil}(F''),\quad A_k\in F'',
\]
retaining all the invisibility assertions of LI.8, and
\[
 f_k(X_3,X_4)=f_k(X_5,X_1)=f_k(X_2,X_3)
       =f_k(X_4,X_5)=f_k(X_1,X_2)=1
       \quad\text{under every }\eta_{p^n}.
\]
Thus its pentagon is one in this entire adelic congruence
family, including every induced five-point quotient, although
$\Delta(k)=1$. This does not assert the full profinite
pentagon for $k$.

\emph{The two harmonic pairs.}
Over $\mathcal O$, the matrices $A,B$ fix the primitive
column $v=(1,1,1)^{\mathsf T}$ and the row
$\ell=(1,\omega,\omega^2)$, with $\ell v=0$.
The basis with columns
\[
 v,\qquad(0,1,-\omega^2)^{\mathsf T},\qquad
                    (0,0,\omega)^{\mathsf T}
\]
has determinant $\omega$, a unit. Hence these matrices are
upper triangular in an integral basis, with diagonal entries
$(1,-\omega,1)$. The group generated by $X_1=A^2,X_2=B^2$
therefore has derived subgroup in the upper unitriangular
group. Over $\mathcal O/p^n$ this is a finite $p$-group.
In particular its second derived group is nilpotent.
The reversing half twist identifies the pair $(X_2,X_3)$
with the first pair, up to interchanging its entries.
Consequently $f\in R_{\rm nil}(F'')$ kills both harmonic
pair evaluations at every prime and every level.

\emph{The three remaining pairs and their uniform seed.}
The following are exact identities over $\mathcal O$:
\[
 (X_4-I)^2=(X_5-I)^2=0,
 \qquad X_i\equiv I\pmod{(1-\omega)}\quad(1\leq i\leq5).
\]
For example,
\[
 X_5=\begin{pmatrix}
 1&\omega-1&-2\omega-1\\0&1&0\\0&0&1
 \end{pmatrix}.
\]
Thus $X_4^{p^n}=X_5^{p^n}=1$ modulo $p^n$.
At $p=3$, the ring $\mathcal O/3^n$ is local with residue
field $\mathbb F_3$ and nilpotent maximal ideal generated
by $1-\omega$. The kernel of its projective linear
reduction to the residue field is a finite three-group.
All the $X_i$ lie in this kernel, so every pair image is a
three-group. Its evaluations are again killed by
$R_{\rm nil}(F'')$.

Now take any finite list of levels $(p,n)$ with $p\ne3$,
and choose a positive integer $m$ divisible by two and by
every listed $p^n$, with $3\nmid m$; their least common
multiple has these properties. Put
\[
             a_m=[[x^m,y^m],[y^m,z^m]]\in F''.
\]
For any group triple with product one, this word is one
if any one of the three displayed $m$th powers is one.
The triples for $(X_3,X_4),(X_5,X_1),(X_4,X_5)$ each
contain $X_4$ or $X_5$. Therefore $a_m$ vanishes at these
three pairs and at all six peripheral permutations, at
every listed level.

These six tests give the full $B$-normal core condition,
not merely a selection of substitutions. Indeed conjugation
by $r$ cyclically permutes $(x,y,z)$, while conjugation
by $s$ swaps $x,y$ and sends $z$ to $(yx)^{-1}$.
The latter three ordered pairs are simultaneously conjugate
to $(y,x),(x,z),(z,y)$, respectively. Thus the six ordered
peripheral pairs represent the six outer cosets in $B/F$;
their kernel intersection is exactly the required $B$-core.

The same word belongs to every full orbifold kernel:
each peripheral permutation of $\varphi$ sends one peripheral
to an element of order two, and $m$ is even. Moreover
\[
                       q_{12}(a_m)=[5I].
\]
This follows directly for the four possible classes
$m\equiv2,4,8,10\pmod{12}$ using
$x\mapsto U^2$, $y\mapsto L_0^2$ and $z=(xy)^{-1}$.
The certificate checks all four classes; equivalently the
modulo-four value is $I$ and the modulo-three value is
$-I$ in each case.

\emph{Actual lifting and compactness.}
Let $K$ be the intersection of the $B$-normal cores of the
three nonharmonic evaluation kernels at the finitely many
listed levels. For any open $B$-normal $M\supseteq F''$
and any finite orbifold target with kernel $J$, set
\[
 N=U_0\cap M\cap K\cap
                   \operatorname{core}_B(\varphi^{-1}J).
\]
This is open and $B$-normal, contains $a_m$, and has
two-group $q_{12}$ image containing $[5I]$.
Apply precisely the coprime $r$-centralizer lift and R2
normalization of LI.8. It supplies an actual normalized
gentle element with the prescribed frame values and
$f,\beta\in R_{\rm nil}(N)$, $A\in N$.
The cusp correction remains in $R_{\rm nil}(N)$, so
it preserves membership in every added evaluation kernel.

Take the compact intersection over all finite lists of
levels, all such $M$, and all $J$. Finite intersections
are nonempty by the construction just given. The LI.8
intersection argument still yields
$f,\beta\in R_{\rm nil}(F'')$, while every nonharmonic
evaluation is one at every level away from three.
The flag argument handles the two harmonic pairs at all
primes, and the residue argument handles every pair at
three. The retained finite frames give $\Delta(k)=1$,
and the full orbifold condition gives $Dk=1$.
The coset representatives used to induce to the full
five-point group may be taken to be powers of the cyclic
rotation of the five $X_i$; hence each block again uses
the same five pairs, all of whose evaluations are one.

For the first odd-field example, an entirely explicit word
which kills even the two harmonic pair maps is
\[
              a_7=\bigl([[x^{28},y^{28}],
                              [y^{28},z^{28}]]\bigr)^7.
\]
It has original value $[5I]$ and all thirty pair/permutation
values one over $\mathbb F_7$. The general proof uses
the uniform unipotent and flag identities, rather than an
enumeration of further finite pentagon assignments.
